% Part: intuitionistic-logic
% Chapter: tableaux
% Section: rules

\documentclass[../../../include/open-logic-section]{subfiles}

\begin{document}

\olfileid{int}{tab}{rul}

\olsection{قواعد المنطق الحدسي}

لا تختلف قواعد~$\land$ و$\lor$ عن قواعد !!{tableau} القضايا المعلّمة
المعهودة إلا بإثبات السوابق فيها. فمتى طُبقت قاعدة على !!{formula}
معلّمة~$\sFmla{\OLId{latin-looped}{S}}{!A}[\sigma]$ حملت كل !!{formula} تنتج منها السابقة
نفسها، أي~$\sigma$. ووجه ذلك ظاهر في التأويل الحدسي: فإن صدقت
$!A \land !B$ في العالم الذي تسميه~$\sigma$، صدقت~$!A$ و$!B$
في~$\sigma$ بعينها لا في غيرها. وقد جُمعت قواعد~$\land$ و$\lor$ في
\olref{tab:prop-rules}.

\begin{table}
  \[\def\arraystretch{3}\begin{array}{|c|c|}
    \hline
    \AxiomC{\sFmla{\True}{!A \land !B}[\sigma]}
    \RightLabel{\TRule{\True}{\land}}
    \UnaryInfC{\sFmla{\True}{!A}[\sigma]}
    \noLine
    \UnaryInfC{\sFmla{\True}{!B}[\sigma]}
    \DisplayProof
    &
    \AxiomC{\sFmla{\False}{!A \land !B}[\sigma]}
    \RightLabel{\TRule{\False}{\land}}
    \UnaryInfC{$\sFmla{\False}{!A}[\sigma] \quad \mid \quad
      \sFmla{\False}{!B}[\sigma]$}
    \DisplayProof
    \\[2ex]
    \hline
    \AxiomC{\sFmla{\True}{!A \lor !B}[\sigma]}
    \RightLabel{\TRule{\True}{\lor}}
    \UnaryInfC{$\sFmla{\True}{!A}[\sigma] \quad \mid \quad
      \sFmla{\True}{!B}[\sigma]$}
    \DisplayProof
    &
    \AxiomC{\sFmla{\False}{!A \lor !B}[\sigma]}
    \RightLabel{\TRule{\False}{\lor}}
    \UnaryInfC{\sFmla{\False}{!A}[\sigma]}
    \noLine
    \UnaryInfC{\sFmla{\False}{!B}[\sigma]}
    \DisplayProof
    \\[2ex]
    \hline
  \end{array}\]
  \caption{قواعد !!{tableau} ذات السوابق للرابطين~$\land$ و$\lor$}
  \ollabel{tab:prop-rules}
\end{table}

وشرط الإغلاق قريب من شرط !!{tableau} العادية، غير أن المطابقة هنا تشمل
السوابق ولا تقتصر على !!{formula}. بل يسوغ لنا شرط أوسع قليلًا: فإن
!!{formula} إذا صدقت في نموذج حدسي دام صدقها في العوالم المتاحة، فلا
يجتمع صدق~$!A$ عند~$\sigma$ وكذبها عند سابقة متاحة~$\sigma.{*}$.
ومن ثم تُغلق الشعبة إذا اجتمع فيها، لبعض
السابقة~$\sigma$ و!!{formula}~$!A$، كلًا من
\[
\sFmla{\True}{!A}[\sigma] \quad\text{و}\quad \sFmla{\False}{!A}[\sigma.{*}]
\]
أما إذا عُكست العلامتان فاشتملت الشعبة على
\[
\sFmla{\False}{!A}[\sigma] \quad\text{و}\quad \sFmla{\True}{!A}[\sigma.{*}]
\]
فلا يوجب ذلك إغلاقها، إلا حين يكون~$*$ هو المتتالية الخالية.

وتُغلق الشعبة كذلك متى احتوت~$\sFmla{\True}{\bot}[\sigma]$.

وأما إدخال الافتراضات فعلى ما هو عليه في !!{tableau} العادية، سوى أننا
نلزمها السابقة~$\OLMathNumeral{1}$؛ ولا خصوصية لهذا العدد، وإنما المقصود اتحاد السابقة
في الافتراضات كلها. ولذلك يكون، مثلًا،
\[
!B_\OLMathNumeral{1}, \dots, !B_\OLId{latin-ordinary}{n} \Proves !A
\]
إذا وفقط إذا وُجد جدول مغلق يبدأ بالافتراضات
\[
\sFmla{\True}{!B_\OLMathNumeral{1}}[\OLMathNumeral{1}], \dots, \sFmla{\True}{!B_\OLId{latin-ordinary}{n}}[\OLMathNumeral{1}],
\sFmla{\False}{!A}[\OLMathNumeral{1}].
\]

وأما قواعد الشرط~$\lif$ فتفارق نظائرها في المنطقين الكلاسيكي
والموجّه. فالقاعدة~$\TRule{\True}{\lif}$ تمد شعبة فيها
$\sFmla{\True}{!A \lif !B}[\sigma]$ بإضافة
$\sFmla{\False}{!A}[\sigma.{*}]$ و$\sFmla{\True}{!B}[\sigma.{*}]$
في شعبتين مختلفتين. ولا تجري إلا بسابقة~$\sigma.{*}$ سبق
\emph{ظهورها} في الشعبة؛ ونسمي مثل هذه السابقة «مستعملة» في تلك
الشعبة. ولما كانت~$\sigma.{*}$ تشمل~$\sigma$ نفسها، أمكن تطبيق القاعدة
دائمًا بإضافة الصيغتين المعلّمتين ذواتي السابقتين
$\sFmla{\False}{!A}[\sigma]$ و$\sFmla{\True}{!B}[\sigma]$
في شعبتين منفصلتين.

والقاعدة~$\TRule{\False}{\lif}$ تمد شعبة فيها
$\sFmla{\False}{!A \lif !B}[\sigma]$ بإضافة كل من
$\sFmla{\True}{!A}[\sigma.\OLId{latin-ordinary}{n}]$ و$\sFmla{\False}{!B}[\sigma.\OLId{latin-ordinary}{n}]$
إلى الشعبة الواحدة، على أن تكون~$\sigma.\OLId{latin-ordinary}{n}$ سابقة لم ترد فيها من قبل.

وقواعد~$\lnot$ نظيرة لهذه، إذ يُؤخذ~$\lnot !A$ بمعنى
$!A \lif \lfalse$.

وهذه القواعد مبسوطة في~\olref{tab:rules-lif-lnot}.

\begin{table}
  \[\def\arraystretch{3}\begin{array}{|c|c|}
    \hline
    \AxiomC{\sFmla{\True}{\lnot !A}[\sigma]}
    \RightLabel{\TRule{\True}{\lnot}}
    \UnaryInfC{$\sFmla{\False}{!A}[\sigma.{*}]$}
    \DisplayProof
    &
    \AxiomC{\sFmla{\False}{\lnot !A}[\sigma]}
    \RightLabel{\TRule{\False}{\lnot}}
    \UnaryInfC{\sFmla{\True}{!A}[\sigma.\OLId{latin-ordinary}{n}]}
    \DisplayProof
    \\[1ex]
    \text{$\sigma.{*}$ مستعملة} & \text{$\sigma.n$ جديدة}\\
    \hline
    \AxiomC{\sFmla{\True}{!A \lif !B}[\sigma]}
    \RightLabel{\TRule{\True}{\lif}}
    \UnaryInfC{$\sFmla{\False}{!A}[\sigma.{*}] \quad \mid \quad
      \sFmla{\True}{!B}[\sigma.{*}]$}
    \DisplayProof
    &
    \AxiomC{\sFmla{\False}{!A \lif !B}[\sigma]}
    \RightLabel{\TRule{\False}{\lif}}
    \UnaryInfC{\sFmla{\True}{!A}[\sigma.\OLId{latin-ordinary}{n}]}
    \noLine
    \UnaryInfC{\sFmla{\False}{!B}[\sigma.\OLId{latin-ordinary}{n}]}
    \DisplayProof
    \\[1ex]
    \text{$\sigma.{*}$ مستعملة} & \text{$\sigma.n$ جديدة}\\
    \hline
  \end{array}\]
  \caption{قواعد !!{tableau} ذات السوابق للرابطين~$\lnot$ و$\lif$}
  \ollabel{tab:rules-lif-lnot}
\end{table}

\end{document}

